\documentclass[11pt,a4paper]{article}
\usepackage[T1]{fontenc}
\usepackage{lmodern,amsmath,amssymb,amsthm,booktabs,array}
\usepackage[margin=1in]{geometry}
\usepackage[colorlinks=true,allcolors=blue]{hyperref}
\usepackage{xurl}
\hypersetup{
pdftitle={Counterexamples to Zoroa et al.'s strong-dominance conjecture for every number of draws above four},
pdfauthor={Haobo Ma; Wenlin Zhang},
pdfsubject={Coupon collection and majorization; CC BY 4.0},
pdfkeywords={coupon collector, stochastic dominance, majorization, counterexample}}
\newtheorem{theorem}{Theorem}
\newtheorem{corollary}[theorem]{Corollary}
\newtheorem{lemma}[theorem]{Lemma}
\newcommand{\Prob}{\mathbb P}
\title{Counterexamples to Zoroa et al.'s strong-dominance conjecture\ for every number of draws above four}
\author{Haobo Ma\\\small ChronoAI Pte Ltd, Singapore
\and Wenlin Zhang\\\small National University of Singapore\\\small ChronoAI Pte Ltd, Singapore}
\date{November 19, 2026}
\begin{document}
\maketitle
\begin{abstract}
We refute the strong-dominance conjecture of Zoroa, Lesigne,
Fern\'andez-S\'aez, Zoroa and Casas \cite[\S8, Conjecture, p.~30,
arXiv:1608.00155v1]{zoroa} for every number of labels $N\geq4$ and every
number of independent draws $n\geq5$, with strictly positive rational
coupon probabilities. For $n=5$ and
$p=(48,48,2,2)/100$ and $q=(50,46,3,1)/100$, one has $p\prec q$, but the
probability mass under $p$ is larger at two distinct labels and smaller at
three distinct labels. These two inequalities exclude every real
single-crossing threshold. The obstruction persists along a rational
one-parameter family on a full real interval. A Taylor estimate extends it
to every number of draws $n\geq5$, and adding sufficiently rare labels
gives strictly positive counterexamples in every dimension $N\geq4$.
The cumulative distributions retain the usual stochastic order. On
supports of at most three distinct counts, that order implies the
conjectured conclusion.
\end{abstract}

\section{The conjecture}
Let $N,n\geq1$ be integers and let $r=(r_1,\ldots,r_N)$ be a probability
vector. Draw $n$ labels independently from $\{1,\ldots,N\}$ with
probabilities $r_i$, and let $Y_n$ count the distinct labels observed.
Write $a_r(k)=\Prob_r(Y_n=k)$ for $k\geq0$.
For two probability vectors, $p\prec q$ means that their decreasing
rearrangements satisfy
\[
 \sum_{i=1}^j p_i^{\downarrow}\leq\sum_{i=1}^j q_i^{\downarrow}
 \quad(1\leq j<N),\qquad
 \sum_{i=1}^N p_i=\sum_{i=1}^N q_i=1.
\]
The definition of strong stochastic dominance in
\cite[Definition~8.2, pp.~28, 30]{zoroa} requires a real $c$ such that
\begin{equation}\label{eq:crossing}
 a_p(k)\leq a_q(k)\quad(k\leq c),\qquad
 a_p(k)\geq a_q(k)\quad(k>c).
\end{equation}
The conjecture states that $p\prec q$ implies \eqref{eq:crossing}.
In the published version these are Definition~7.2 and the conjecture
in \S7; the version-specific locators above refer to the arXiv source.
Here ``strong'' refers to this single-crossing condition, not merely to
the usual stochastic order
$\Prob_p(Y_n\leq k)\leq\Prob_q(Y_n\leq k)$.
The latter was already proved by Wong and Yue
\cite[pp.~391--398]{wongyue}, as explicitly credited in
\cite[Remark~6.1, pp.~25--26]{zoroa}; see also
\cite[Corollary~6.1, p.~25]{zoroa}.
For related pairwise-mixing inequalities for collection times, see
\cite[Theorem~1]{anceaume}. Lesigne's later exposition
\cite[\S4, p.~5; \S7, Proposition~1]{lesigne} presents the cumulative
comparison and the inclusion--exclusion mass law. The distinction matters: an order of
cumulative probabilities need not give an order of the individual masses
on each side of one threshold.

\section{Four labels and five draws}
\begin{theorem}\label{thm:counterexample}
For $N=4$, $n=5$ and
\[
 p=\frac1{100}(48,48,2,2),\qquad
 q=\frac1{100}(50,46,3,1),
\]
all coordinates are positive and $p\prec q$, but
\eqref{eq:crossing} fails for every $c\in\mathbb R$.
\end{theorem}
\begin{proof}
Both vectors are decreasing and their coordinate sums are one. Their
prefix sums, multiplied by $100$, are respectively
\[
 (48,96,98,100),\qquad(50,96,99,100),
\]
so $p\prec q$.

We compute the masses without approximation. For $r\in\mathbb R^4$ with
nonnegative coordinates summing to one, put
\[
 S_j(r)=\sum_{\substack{A\subseteq\{1,2,3,4\}\\|A|=j}}
              \left(\sum_{i\in A}r_i\right)^5
 \qquad(1\leq j\leq3).
\]
The event that the set of observed labels is exactly a fixed set $A$
has probability
\begin{equation}\label{eq:ie}
 \sum_{B\subseteq A}(-1)^{|A|-|B|}
                 \left(\sum_{i\in B}r_i\right)^5.
\end{equation}
Indeed, the fifth power is the probability that every draw lies in $B$;
inclusion--exclusion removes the events in which a label of $A$ is absent.
The empty-set term is zero because there are five draws. Summing
\eqref{eq:ie} over sets $A$ of a fixed size gives
\begin{align*}
 a_r(1)&=S_1(r),\\
 a_r(2)&=S_2(r)-3S_1(r),\\
 a_r(3)&=S_3(r)-2S_2(r)+3S_1(r),\\
 a_r(4)&=1-a_r(1)-a_r(2)-a_r(3).
\end{align*}
For example, each singleton lies in three pairs, each pair in two triples,
and each singleton in three triples; these counts give the coefficients.
The required integer sums are
\[
\begin{array}{c|rrr}
 &100^5S_1&100^5S_2&100^5S_3\\\hline
 p&509\,608\,000&9\,403\,728\,000&18\,838\,824\,000\\
 q&518\,463\,220&9\,428\,769\,000&18\,868\,905\,780
\end{array}
\]
and hence the full mass functions have the following numerators over the
common denominator $100^5=10^{10}$:
\begin{center}
\begin{tabular}{crrr}
\toprule
$k$ & $100^5a_p(k)$ & $100^5a_q(k)$ & $100^5(a_p(k)-a_q(k))$\\
\midrule
1 & $509\,608\,000$ & $518\,463\,220$ & $-8\,855\,220$\\
2 & $7\,874\,904\,000$ & $7\,873\,379\,340$ & $1\,524\,660$\\
3 & $1\,560\,192\,000$ & $1\,566\,757\,440$ & $-6\,565\,440$\\
4 & $55\,296\,000$ & $41\,400\,000$ & $13\,896\,000$\\
\bottomrule
\end{tabular}
\end{center}
For $k=0$ or $k>4$ both masses vanish. In particular, $a_p(2)>a_q(2)$.
If a threshold $c$ existed, the lower inequality in
\eqref{eq:crossing} would force $c<2$. Then $3>c$ would require
$a_p(3)\geq a_q(3)$, contrary to the third row. This proves the theorem.
\end{proof}

The cumulative differences, multiplied by $100^5$, are
$-8\,855\,220$, $-7\,330\,560$, $-13\,896\,000$, and $0$ at
$k=1,2,3,4$. Thus this example retains ordinary stochastic dominance;
it separates that order from the source's stronger single-crossing condition.

\begin{corollary}\label{cor:family}
For every real $0<t\leq1/100$, the positive probability vectors
\[
 p_t=(1/2-2t,1/2-2t,2t,2t),\qquad
 q_t=(1/2,1/2-4t,3t,t)
\]
satisfy $p_t\prec q_t$ and fail \eqref{eq:crossing} after five draws.
\end{corollary}
\begin{proof}
Both vectors are decreasing on this interval, and their prefix sums are
\[
 (1/2-2t,1-4t,1-2t,1),\qquad(1/2,1-4t,1-t,1).
\]
This proves majorization; all coordinates are positive.
Substituting in the formulas for $S_j$ gives the exact differences
\begin{align*}
 a_{p_t}(2)-a_{q_t}(2)&=t^2(5-360t+1270t^2-2340t^3),\\
 a_{p_t}(3)-a_{q_t}(3)&=t^2(-10+360t-1670t^2+1560t^3).
\end{align*}
The first bracket is at least
$5-360/100-2340/100^3=69883/50000>0$.
The second is at most
$-10+360/100+1560/100^3=-159961/25000<0$.
As $t^2>0$, the same threshold contradiction applies. At $t=1/100$ these
vectors equal those of Theorem~\ref{thm:counterexample}; at the excluded
endpoint $t=0$ they coincide.
\end{proof}

\section{Every number of draws above four}
The family in Corollary~\ref{cor:family} also explains why the obstruction
is not confined to five draws. In this section $a_r(k)$ denotes the mass
after $n$ draws, and $S_j(r;n)$ is the subset sum in the preceding proof
with exponent $n$ in place of $5$.

\begin{theorem}\label{thm:alln}
For every integer $n\geq5$ and every real
$0<t\leq (100n2^n)^{-1}$, the strictly positive four-coordinate vectors
$p_t,q_t$ of Corollary~\ref{cor:family} satisfy $p_t\prec q_t$ and
\[
 a_{p_t}(2)-a_{q_t}(2)>0,\qquad
 a_{p_t}(3)-a_{q_t}(3)<0.
\]
In particular they fail the single-crossing condition after $n$ draws.
\end{theorem}
\begin{proof}
The prefix inequalities in Corollary~\ref{cor:family} hold for
$0<t\leq1/14$, which contains the interval in the theorem.
Put $G_j(t)=S_j(p_t;n)-S_j(q_t;n)$ and
\[
 F_2(t)=G_2(t)-3G_1(t),\qquad
 F_3(t)=G_3(t)-2G_2(t)+3G_1(t).
\]
Inclusion--exclusion identifies these with the two mass differences.
Every term of $G_j$ is a signed power $(b+st)^n$, where the base $b$
is $0$, $1/2$ or $1$. The constant and linear coefficients of each
$G_j$ vanish. Write $[t^2]G_j$ for the coefficient of $t^2$ in $G_j$.
With $C=\binom n2$ and $x=2^{2-n}$, their quadratic
coefficients and the sums of absolute cubes of their slopes are
\[
\begin{array}{c|rr}
 & [t^2]G_j & \sum |s|^3\\\hline
G_1&-8Cx&124\\
G_2&-20Cx&312\\
G_3&C(-8x-2)&124
\end{array}
\]
Here the last column includes terms from both vectors. To see the first
row, the slopes on the two coordinates of base $1/2$ are $-2,-2$
for $p_t$ and $0,-4$ for $q_t$, whose squared-slope difference is $-8$.
For pairs of base $1/2$, the four slopes are all zero for $p_t$, and
$3,1,-1,-3$ for $q_t$, giving $-20$. For triples the squared-slope
differences are $-8$ at base $1/2$ and $-2$ at base $1$.
Terms of base zero have no quadratic contribution since $n\geq5$.
Thus
\begin{align}
 F_2(t)&=n(n-1)2^{3-n}t^2+R_2(t),\label{eq:alln2}\\
 F_3(t)&=n(n-1)(2^{4-n}-1)t^2+R_3(t).\label{eq:alln3}
\end{align}
All subset probabilities lie in $[0,1]$ for $0\leq t\leq1/14$.
The third derivative of $(b+st)^n$ therefore has absolute value at
most $n(n-1)(n-2)|s|^3$. Taylor's theorem gives
\[
 |R_2(t)|\leq114n(n-1)(n-2)t^3,\qquad
 |R_3(t)|\leq\frac{560}{3}n(n-1)(n-2)t^3,
\]
because $(312+3\cdot124)/6=114$ and
$(124+2\cdot312+3\cdot124)/6=560/3$.
Since $(n-2)t\leq 2^{-n}/100$ and $n\geq5$, these estimates yield
\begin{align*}
 F_2(t)&\geq \frac{343}{50}n(n-1)2^{-n}t^2>0,\\
 F_3(t)&\leq n(n-1)t^2\left(-\frac12+\frac7{120}\right)
          =-\frac{53}{120}n(n-1)t^2<0.
\end{align*}
The positive difference at two forces a proposed threshold below two,
and the negative difference at three then contradicts its upper side.
\end{proof}

\begin{corollary}\label{cor:allN}
For every pair of integers $N\geq4$ and $n\geq5$, there are strictly
positive rational probability vectors $p\prec q$ on $N$ labels for which
the single-crossing condition fails after $n$ draws.
\end{corollary}
\begin{proof}
For $N=4$, take $t=(100n2^n)^{-1}$ in Theorem~\ref{thm:alln}.
For $M=N-4>0$, use the same $t$ and put
\[
 m=n(n-1)2^{-n}t^2,\qquad
 \delta=\min\{m/(4n),Mt/4,1/4\}.
\]
Multiply the four coordinates of each vector by $1-\delta$ and
append $M$ coordinates equal to $\delta/M$. All coordinates are
positive rational numbers. Since $\delta/M\leq t/4$ and
$(1-\delta)t\geq3t/4$, the appended coordinates are no larger than
the original four. The first four prefix comparisons are just the
old comparisons scaled by $1-\delta$; all subsequent prefixes agree.
Hence majorization is preserved.

Couple a draw from the enlarged vector with a draw from its original
four-coordinate vector: with probability $1-\delta$ use the same
original label, and otherwise use one of the appended labels.
For $n$ draws the probability that any appended label occurs is at
most $n\delta$. Consequently each mass changes by at most $n\delta$,
and each difference of two masses changes by at most $2n\delta\leq m/2$.
The estimates in Theorem~\ref{thm:alln} give $F_2(t)\geq m$ and
$F_3(t)\leq-m$, so their opposite signs persist.
\end{proof}

\section{Small supports and coverage}
\begin{lemma}\label{lem:three}
For distributions supported on $\{1,2,3\}$, ordinary stochastic
dominance implies the single-crossing condition \eqref{eq:crossing}.
Consequently, for coupon collection, the conjectured conclusion holds
whenever $\min(N,n)\leq3$ and $p\prec q$.
\end{lemma}
\begin{proof}
Let $d_k=a_p(k)-a_q(k)$. The cumulative inequalities give
$d_1\leq0$ and $d_1+d_2\leq0$; equality of total masses then gives
$d_3=-(d_1+d_2)\geq0$.
If $d_2\geq0$, choose $c=1$; otherwise choose $c=2$.
In either case \eqref{eq:crossing} follows, including zero masses
outside $\{1,2,3\}$.
For coupon collection the support is contained in
$\{1,\ldots,\min(N,n)\}$, and ordinary stochastic dominance follows
from Wong and Yue \cite[pp.~391--398]{wongyue}, also restated in
\cite[Corollary~6.1 and Remark~6.1, pp.~25--26]{zoroa}.
\end{proof}
Thus at least four labels and four possible distinct counts are needed
for a counterexample. Theorem~\ref{thm:alln} covers every $n\geq5$;
the case $n=4$, $N>5$ is not determined here.

\begin{center}
\begin{tabular}{p{.47\linewidth}p{.44\linewidth}}
\toprule
Source item in \cite{zoroa} & Coverage\\
\midrule
\S8, Conjecture, p.~30 & Refuted (explicit counterexamples),
Theorems~\ref{thm:counterexample}, \ref{thm:alln} and Corollary~\ref{cor:allN}.\\
\S8, $n=2,3$ for arbitrary $N$, p.~30 &
Reproved; prior proof Zoroa et al. (2016), deduced in Lemma~\ref{lem:three}.\\
\S8, $n=4$, $N\leq5$, p.~30 &
Reported there as proved, without a proof in that passage;
not covered here beyond $N\leq3$.\\
Corollary~6.1, p.~25 &
Prior proof: Wong and Yue (1973) \cite[pp.~391--398]{wongyue},
credited in \cite[Remark~6.1, pp.~25--26]{zoroa};
ordinary stochastic order retained.\\
Extension to $n\geq5$, $N\geq4$ &
Not posed; proved in Theorem~\ref{thm:alln} and Corollary~\ref{cor:allN}.\\
\bottomrule
\end{tabular}
\end{center}

\paragraph{Verification coverage.}
The declaration \texttt{result : $\neg$ claim} in the
\href{https://github.com/the-omega-institute/trureturing/blob/ad59550e53d607f9e9dbd8110a79414c46e97f7e/D5/S3/Analytic/Probability/StrongDominanceRefutation.lean}{pinned Lean source}
proves \texttt{not claim} from the vectors of
Theorem~\ref{thm:counterexample} at $N=4$, $n=5$ only.
It formalizes their decreasing-prefix inequalities and the four mass
values $a_p(2)$, $a_q(2)$, $a_p(3)$ and $a_q(3)$.
The masses at $k=1,4$, the displayed cumulative differences,
Corollaries~\ref{cor:family} and \ref{cor:allN}, and
Theorem~\ref{thm:alln} have written proofs and script checks only;
they are outside this declaration's scope.
The inclusion--exclusion formulas and Lemma~\ref{lem:three} have written
proofs. The accompanying \texttt{verify.py} checks the concrete
mass table, the exact family polynomials and interval bounds, and the
coefficient weights and derivative constants in the all-$n$ proof.
Finite endpoint checks are distinguished from the all-$n$ proof in
\texttt{verification.json}.

\begin{thebibliography}{9}
\bibitem{wongyue}
C.~K.~Wong and P.~C.~Yue,
\emph{A majorization theorem for the number of distinct outcomes in
$n$ independent trials}, Discrete Mathematics \textbf{6} (1973),
391--398.
\href{https://doi.org/10.1016/0012-365X(73)90070-8}{doi:10.1016/0012-365X(73)90070-8}.
\bibitem{zoroa}
N.~Zoroa, E.~Lesigne, M.~J.~Fern\'andez-S\'aez, P.~Zoroa and J.~Casas,
\emph{The coupon collector urn model with unequal probabilities in ecology
and evolution}, Journal of the Royal Society Interface \textbf{14} (2017),
20160643. \href{https://doi.org/10.1098/rsif.2016.0643}{doi:10.1098/rsif.2016.0643}.
Posing version: \href{https://arxiv.org/abs/1608.00155v1}{arXiv:1608.00155v1}
(2016), \S8, pp.~28--30.
\bibitem{anceaume}
E.~Anceaume, Y.~Busnel, E.~Schulte-Geers and B.~Sericola,
\emph{Optimization results for a generalized coupon collector problem},
Journal of Applied Probability \textbf{53} (2016), 622--629.
\href{https://doi.org/10.1017/jpr.2016.27}{doi:10.1017/jpr.2016.27}.
Author version: \href{https://arxiv.org/abs/1504.03878v1}{arXiv:1504.03878v1}, Theorem~1.
\bibitem{lesigne}
E.~Lesigne, \emph{Collection de vignettes in\'egalement distribu\'ees ou
quand l'h\'et\'erog\'en\'eit\'e est une protection},
\href{https://hal.science/hal-05746525v1}{HAL:hal-05746525v1} (2026),
\S4, p.~5 and \S7, pp.~9--11.
\end{thebibliography}
\noindent\footnotesize Copyright \copyright\ 2026 Haobo Ma and Wenlin Zhang. License: CC BY 4.0.
\end{document}
